\documentclass[10pt,a4paper]{amsart}
\usepackage[margin=19mm]{geometry}
\usepackage{amsmath,amssymb,mathtools}
\usepackage[scr=rsfs,cal=euler]{mathalfa}
\usepackage{xfrac}
\usepackage[unicode,hidelinks,bookmarks=false]{hyperref}
\newcommand{\ie}{i.e.\ }
\newcommand{\cf}{cf.\ }
\newcommand{\resp}{respectively\ }
\newcommand{\Subsection}{\subsection}
\providecommand{\nobreakdash}{\nobreak\hbox{-}}
\setlength{\emergencystretch}{2em}
\multlinegap0pt
\usepackage{bm}
\usepackage{bbm}
\usepackage{dsfont}
\usepackage{xspace}
\usepackage{cancel}
\usepackage{mathtools}
\usepackage{amsthm}
\makeatletter
\@ifundefined{thm}{\newtheorem{thm}{Thm}}{}
\@ifundefined{lemm}{\newtheorem{lemm}[thm]{Lemma}}{}
\makeatother
\renewcommand{\setminus}{\smallsetminus}
\title[Inner Lipschitz Geometry of Complex Surface Germs with Non-I]{Inner Lipschitz Geometry of Complex Surface Germs with Non-Isolated Singularities: A Complete Classification}
\author{Yenni Cherik}
\date{Source: arXiv:2606.15941v1 (CC BY 4.0)}
\begin{document}
\maketitle

\noindent\emph{Source and licence.} Inner Lipschitz Geometry of Complex Surface Germs with Non-Isolated Singularities: A Complete Classification. Version arXiv:2606.15941v1 (CC BY 4.0), distributed under the Creative Commons Attribution 4.0 International licence, as recorded in the arXiv licence metadata for this exact version and shown on the abstract page https://arxiv.org/abs/2606.15941v1. Full rights notice: licenses/arxiv-2606.15941-CC-BY-4.0.txt.\par\smallskip

\noindent\textit{Editorial scope.} Counted unit: the Lemm whose source label is \texttt{keylemma} in the article (source0.tex lines 4290--4309, source label \texttt{keylemma}), delivered together with the article's own complete original proof of it (source0.tex lines 4311--4415, one \texttt{proof} environment of 105 lines). No mathematics is written, repaired or re-derived: statement and proof are the article's own text line for line. Required context retained uncounted: none. Every definition, display and label of those lines is retained unchanged. Omitted from this package and not redistributed: 1--4289, 4310--4310, 4416--4723. Nothing inside a retained excerpt is omitted or altered: every retained body is delivered exactly as the article's lines are written, so no omission receipt of any kind accompanies this package. Citations: the retained excerpts carry no citation command at all, so no bibliography entry of the article is transcribed and no reference is redistributed. Reference adaptation: none was needed and none was made; every reference inside the delivered unit resolves inside it, so no unresolved reference or citation is produced. Printed slips kept literal: no printed word, symbol, hypothesis or number of the counted statement or of its proof was altered. Layout and class: the article's own class is \texttt{amsart} with packages including geometry, amsmath, amssymb, mathrsfs, amscd, subcaption, ulem, babel, fancyhdr, lmodern, float, stmaryrd, hyperref, enumitem; the wrapper is the lane's pinned \texttt{amsart} editorial wrapper at 10pt on A4, which restates the theorem-like environments the retained text needs and adds the article's own macro definitions that the retained text needs (\texttt{\textbackslash setminus}), with extra wrapper packages (bm, bbm, dsfont, xspace, cancel, mathtools). The delivered numbering is the unit's own. Assets: no retained span contains a figure environment, a graphics-inclusion command, a PDF-page inclusion, a table or a TikZ picture; no image file is copied into this package, the package declares no asset dependency and no image file is redistributed with it. Licence: version arXiv:2606.15941v1 of this article is distributed under the Creative Commons Attribution 4.0 International licence, as recorded in the arXiv licence metadata for this exact version and shown on the abstract page https://arxiv.org/abs/2606.15941v1. A full rights notice is delivered at licenses/arxiv-2606.15941-CC-BY-4.0.txt. Characters: every retained excerpt is pure ASCII, so the delivered engine, XeLaTeX with its default Latin Modern text font, reports no missing character.
\begin{lemm}\label{keylemma}
Let $(X,0)$ be a complex surface germ and let 
\begin{center}
\begin{tabular}{ccc}
$n:(\overline{X},0)$ & $\longrightarrow$ & $(X,0) \subset (\mathbb{C}^k,0)$ \\
$p$ & $\longmapsto$ & $(f_1(p),f_2(p),\dots,f_k(p))$
\end{tabular}
\end{center}
be its normalization. Then the following are equivalent:
\begin{enumerate}
  \item The generating system $(f_1,\dots,f_k)$ is precomplete for the ideal $I:=n^*(\mathfrak{m})$, where $\mathfrak{m}$ is the maximal ideal of $\mathcal{O}_{X,0}$. \label{ma}
    \item The map $n$
    is an immersion outside $\{0\}$.\label{ho}
    
  
    
    \item The family $(\mathrm{d}f_1,\dots,\mathrm{d}f_k)$ generates the $\mathcal{O}_{X,0}$-submodule $\Omega_I^1 := \{ \mathrm{d}h \mid h \in I \}
    \subset \Omega^1_{\overline{X},0}.$\label{gani}
\end{enumerate}
\end{lemm}

\begin{proof}


\medskip
\noindent\emph{(i)$\Rightarrow$(ii).}
 We want to prove that the map $n$ is an immersion on $\overline{X} \backslash \{0\}$. By contradiction, suppose that $n$ is not an immersion. Then there exists $p \in \overline{X} \setminus {0}$ such that
\[
\mathrm{d}f_i \wedge \mathrm{d}f_j(p) = 0 \quad \text{for all } i,j = 1,\dots,k.
\]

Let $(\ell_1,\dots,\ell_r)$ be a system of generators of the maximal ideal $\mathfrak{m}$ of $\mathcal{O}_{X,0}$ such that the morphism
\[
\begin{array}{rcl}
\ell: (\overline{X} ,0) &\to& (\mathbb{C}^r,0) \\
p &\mapsto &(\ell_1(p),\dots,\ell_r(p))
\end{array}
\]
is an embedding. In particular, there exist indices $i,j,m \in \{1,\dots,r\}$ such that $\ell_m(p) \neq 0$ and
\[
\mathrm{d}\ell_i \wedge \mathrm{d}\ell_j(p) \neq 0.
\]

By hypothesis, there exists $n \in \mathbb{N}$ such that $\mathfrak{m}^n \subset I$. Consider the functions
\[
F_{im}^{(ab)} = (a\ell_i + b\ell_m)^n, \qquad F_{jm}^{(cd)} = (c\ell_j + d\ell_m)^n,
\]
with $a,b,c,d \in \mathbb{C}$. Since $F_{im}^{(ab)}$ and $F_{jm}^{(cd)}$ are elements of $I$, their wedge product of differentials
\[
\omega = \mathrm{d}F_{im}^{(ab)} \wedge \mathrm{d}F_{jm}^{(cd)}
\]
belongs to $\Omega_I^2$, and explicitly,
\[
\omega = n^2 (a\ell_i + b \ell_m)^{n-1} (c\ell_j + d \ell_m)^{n-1} 
\bigl( ac \, \mathrm{d}\ell_i \wedge \mathrm{d}\ell_j + ad \, \mathrm{d}\ell_i \wedge \mathrm{d}\ell_m + bc \, \mathrm{d}\ell_m \wedge \mathrm{d}\ell_j \bigr).
\]

We can choose $a,b,c,d$ such that $\omega(p) \neq 0$. By assumption, $(f_1,\dots,f_k)$ is a precomplete system of generators of $I$, so there exists a family of functions $\{\psi_{ij}\}_{1\le i,j \le k}$ such that
\[
\omega = \sum_{1 \le i,j \le k} \psi_{ij} \, \mathrm{d}f_i \wedge \mathrm{d}f_j.
\]

However, this equality cannot hold at $p$ if $\mathrm{d}f_i \wedge \mathrm{d}f_j(p) = 0$ for all $i,j = 1,\dots,k$. This contradiction shows that $n$ must indeed be an immersion on $\overline{X} \backslash \{0\}$.


\medskip
\noindent\emph{(ii)$\Rightarrow$(iii).} Assume $n : (\overline{X},0)\to (X,0)\subset (\mathbb{C}^k,0)$ is an immersion on $\overline{X}\setminus{0}$. It is clear that $(\mathrm{d}f_1,\dots,\mathrm{d}f_k)\subset \Omega^1_I$, so it remains to prove the reverse inclusion.
Since $n$ is a normalization, it is bimeromorphic. In particular, the induced morphism
$$n^* : \mathcal{M}(X,0) \to \mathcal{M}(\overline{X},0), n^*(g):=g \circ n$$
is an isomorphism, where $\mathcal{M}(X,0)$ (resp. $\mathcal{M}(\overline{X},0)$) denotes the field of meromorphic function germs.
Let $h \in I \subset \mathcal{O}_{\overline{X},0} \subset \mathcal{M}(\overline{X},0)$. Then there exists a unique $g \in \mathcal{M}(X,0)$ such that
'$h = g \circ n$.
Thus,
$h = \sum_{i=1}^k \frac{\partial g}{\partial z_i}(f_1,\dots,f_k) \mathrm{d}f_i$,
with $\frac{\partial g}{\partial z_i} \in \mathcal{M}(X,0)$. Hence the coefficients
$c_i := \frac{\partial g}{\partial z_i}(f_1,\dots,f_k)$
belong to $\mathcal{M}(\overline X,0)$.
On the other hand, since $h \in I = (f_1,\dots,f_k)$, we can write
$h = \sum_i a_i f_i$, with $a_i \in \mathcal{O}_{\overline X,0}$, and therefore
$\mathrm{d}h = \sum_i a_i \mathrm{d}f_i + \sum_i f_i \mathrm{d}a_i$.
Let $C \subset \overline{X}$ be an irreducible complex curve through $0$ not contained in the singular locus. Restricting the above identities to $C$, we obtain
$\mathrm{d}h|_C = \sum_i (c_i|_C) \mathrm{d}f_i|_C$.
Since $n$ is an immersion on $\overline X \setminus {0}$, there exists $i$ such that $\mathrm{d}f_i|_C \not\equiv 0$. As $\mathrm{d}h|_C$ is holomorphic at $0$, it follows that $c_i|_C$ has no pole at $0$. Hence each $c_i$ is holomorphic along every such curve $C$.
Since $\overline{X}$ is normal, this implies that $c_i \in \mathcal{O}_{\overline{X},0}$. Therefore,
$\mathrm{d}h = \sum_i c_i, \mathrm{d}f_i$
with $c_i \in \mathcal{O}_{\overline{X},0}$, and thus $\Omega^1_I \subset (\mathrm{d}f_1,\dots,\mathrm{d}f_k)$.







%Now, let us prove that $\Omega_I^1 \subset M$. For any \( N \in \mathbb{N} \), any element \( g \in I \) can be written as
%\[
%g = \sum_i a_i f_i + h, \quad a_i \in \mathbb{C},\ h \in \mathfrak{m}^{N+1} I.
%\]

%Differentiating, we have
%\[
%dg = \sum_i a_i\, df_i + %dh.
%\]

%Write $h = \ell \widetilde{h}$ with $\ell \in \mathfrak{m}^{N+1}$ and $\widetilde{h} \in I$. Then
%\[
%dh = \ell\, d\widetilde{h} + \widetilde{h}\, d\ell.
%\]

%The first term satisfies $\ell\, d\widetilde{h} \in \mathfrak{m}^{N+1} \Omega^1_{X,0} \subset M$, and the second term $\widetilde{h}\, d\ell \in I\, \mathrm{d}(\mathfrak{m}^{N+1}) \subset I \mathfrak{m}^N \Omega^1_{X,0} \subset M$. Therefore $dg \in M$, so $\Omega_I^1 \subset M$.















\noindent\emph{(iii)$\Rightarrow$(i).} It follows immediately from the definition of precomplete generator system for an ideal.
\end{proof}

\end{document}
